\section{One-Time Pad (OTP)}
\label{sec:otp}

\subsection{Adakah Cipher yang Tidak Dapat Dipecahkan?}

\crypt{Unbreakable cipher} adalah klaim yang sering dibuat kriptografer terhadap algoritma rancangannya. Kenyataannya, hampir semua algoritma yang pernah dibuat adalah \textit{breakable cipher}: Caesar, Vigen\`{e}re, Playfair, Affine, Hill, dan Enigma yang dibahas pada bab ini semuanya sudah dapat dipecahkan \cite{munir_slides_itb}.

Apakah \textit{unbreakable cipher} benar-benar ada? Jawabannya \textbf{ada}, asalkan tiga syarat berikut dipenuhi:
\begin{enumerate}
  \item kunci dibangkitkan secara \textbf{acak sejati} (\textit{truly random});
  \item \textbf{panjang kunci sama dengan panjang plainteks};
  \item kunci \textbf{hanya dipakai sekali} dan tidak boleh dipakai ulang untuk mengenkripsi pesan lain.
\end{enumerate}

Acak sejati berarti nilainya tidak dapat diprediksi dan pembangkitannya tidak dapat diulang. Ini berbeda dengan bilangan acak semu (\textit{pseudo-random}), yang deretnya dapat dibangkitkan ulang persis sama asalkan umpannya (\textit{seed}) sama. Akibat syarat 1 dan 2, plainteks yang sama tidak selalu menghasilkan cipherteks yang sama.

\subsection{Sejarah: Vernam dan Shannon}

Satu-satunya algoritma kriptografi yang aman sempurna (\textit{perfect secrecy}) sehingga tidak dapat dipecahkan adalah \crypt{one-time pad} (OTP). OTP ditemukan oleh \textbf{Gilbert Vernam} (Gambar~\ref{fig:potret-vernam}) pada tahun 1917, lalu dibuktikan secara matematis memiliki \textit{perfect secrecy} oleh \textbf{Claude Shannon} pada tahun 1949 \cite{munir_slides_itb,shannon1949}.

\begin{figure}[htbp]
  \centering
  \includegraphics[width=0.2\textwidth]{bab-04/potret-vernam}
  \caption{Gilbert Sandford Vernam (1890--1960), insinyur AT\&T Bell Labs yang pada 1917 merancang cipher teleprinter dengan kunci pada pita kertas \cite{munir_slides_itb}.}
  \label{fig:potret-vernam}
\end{figure}

Vernam merancang cipher teleprinter: kunci yang disiapkan sebelumnya disimpan pada pita kertas, lalu digabungkan karakter demi karakter dengan plainteks untuk menghasilkan cipherteks. Untuk dekripsi, kunci yang sama digabungkan kembali dengan cipherteks.

OTP mengatasi kelemahan Vigen\`{e}re cipher (Subbab~\ref{sec:vigenere}). Vigen\`{e}re mengulang kunci secara periodik sehingga panjang kuncinya mudah ditemukan dengan metode Kasiski. Pada OTP, panjang kunci sama dengan panjang plainteks dan seluruh hurufnya acak sejati, misalnya:
\begin{center}
\ttfamily
\begin{tabular}{ll}
  \rmfamily Plainteks: & otpadalahcipheryangtidakbisadipecahkan \\
  \rmfamily Kunci: & trjkdndkdwerylgrgdkopcegyhbdwjbtrfhgvk
\end{tabular}
\end{center}

\subsection{Pad: Bloknot Kunci Sekali Pakai}

\textit{Pad} berarti kertas bloknot. Sebuah \textit{one-time pad} berisi deretan huruf kunci yang dibangkitkan secara acak (Gambar~\ref{fig:otp-pad}), misalnya:
\begin{center}
\small\ttfamily
\begin{tabular}{l}
CIHJT UUHML FRUGC ZIBGD BQPNI PDNJG LPLLP YJYXM \\
DCXAC JSJUK BIOYT MWQPX DLIRC BEXYK VKIMB TYIPE \\
UOLYQ OKOXH PIJKY DRDBC GEFZG UACKD RARCD HBYRI
\end{tabular}
\end{center}

\begin{figure}[htbp]
  \centering
  \includegraphics[width=0.6\textwidth]{bab-04/otp-pad-cryptomuseum}
  \caption{Bloknot \textit{one-time pad} berisi kelompok-kelompok huruf kunci acak (sumber: \url{https://www.cryptomuseum.com/crypto/otp/index.htm}) \cite{munir_slides_itb}.}
  \label{fig:otp-pad}
\end{figure}

Cara pemakaiannya \cite{munir_slides_itb}:
\begin{itemize}
  \item pengirim dan penerima pesan memiliki salinan \textit{pad} yang sama;
  \item satu lembar \textit{pad} hanya dipakai sekali (\textit{one-time}) untuk mengenkripsi pesan, dari sinilah nama \textit{one-time pad} berasal;
  \item setelah dipakai, lembar \textit{pad} dihancurkan agar tidak dipakai lagi untuk pesan lain, sehingga kriptanalisis menjadi sulit.
\end{itemize}

\subsection{Enkripsi dan Dekripsi}

Aturan enkripsi dan dekripsi OTP persis sama dengan Vigen\`{e}re cipher; bedanya, tidak ada perulangan kunci secara periodik. Dengan pengodean huruf pada Tabel~\ref{tab:kode-huruf}:
\[
\text{Enkripsi: } C = (P + K) \bmod 26 \qquad
\text{Dekripsi: } P = (C - K) \bmod 26
\]
Contoh: untuk $P = \texttt{o}$ dan $K = \texttt{t}$, $C = (14 + 19) \bmod 26 = 33 \bmod 26 = 7 = \texttt{H}$. Sebaliknya, $P = (7 - 19) \bmod 26 = -12 \bmod 26 = 14 = \texttt{o}$.

Untuk data biner, penjumlahan modulo 26 diganti dengan XOR bit demi bit \cite{katz2020,stallings2017}:
\[
c_i = p_i \oplus k_i, \qquad p_i = c_i \oplus k_i.
\]
Bentuk inilah yang menjadi model ideal \textit{stream cipher} pada Bab~V.

\begin{contoh}
\textbf{Contoh 1.} Plainteks \texttt{onetimepad}, kunci \texttt{tbfrgfarfm}:
\begin{align*}
(\texttt{o} + \texttt{t}) \bmod 26 &= (14 + 19) \bmod 26 = 7 = \texttt{H} \\
(\texttt{n} + \texttt{b}) \bmod 26 &= (13 + 1) \bmod 26 = 14 = \texttt{O} \\
(\texttt{e} + \texttt{f}) \bmod 26 &= (4 + 5) \bmod 26 = 9 = \texttt{J}, \quad \text{dan seterusnya.}
\end{align*}
Cipherteks: \texttt{HOJKOREGFP}.

\medskip
\textbf{Contoh 2.} Pesan yang lebih panjang:
\begin{center}
\small\ttfamily
\begin{tabular}{ll}
  \rmfamily Plainteks: & nantimalamsayatunggukamudidepanwarungkopi \\
  \rmfamily Kunci: & gtrskncvbrwpoatqljfmxrpjsrzfoltfbtmaedpvy \\
  \rmfamily Cipherteks: & TTELSZCGBDOPMAMKYPLGHRBDVZCJDLGBBKGNKNDKG
\end{tabular}
\end{center}
\end{contoh}

\begin{catatan}
\textbf{Kakas daring.} Enkripsi dan dekripsi OTP dapat dicoba di \url{https://www.cachesleuth.com/onetimepad.html} (Gambar~\ref{fig:otp-demo}). Kakas ini juga dapat membangkitkan kunci acak sepanjang pesan.
\end{catatan}

\begin{figure}[htbp]
  \centering
  \includegraphics[width=0.75\textwidth]{bab-04/otp-demo-cachesleuth}
  \caption{Kakas daring \textit{One Time Pad Cipher} dari CacheSleuth \cite{munir_slides_itb}.}
  \label{fig:otp-demo}
\end{figure}

\subsection{Mengapa OTP Tidak Dapat Dipecahkan?}

Ada dua alasan \cite{munir_slides_itb}:
\begin{enumerate}
  \item \textbf{Kunci acak + plainteks tidak acak = cipherteks acak.} Statistik bahasa pada plainteks tertutup sepenuhnya oleh kunci yang acak.
  \item \textbf{Hanya ada satu kunci yang memetakan sebuah plainteks ke sebuah cipherteks}, dan sebaliknya. Fungsi $C = (P + K) \bmod 26$ adalah pemetaan satu-ke-satu. Akibatnya, jika kriptanalis mencoba mendekripsi cipherteks dengan kunci-kunci yang berbeda, ia bisa memperoleh beberapa plainteks yang sama-sama bermakna dan tidak dapat menentukan mana yang benar.
\end{enumerate}

\begin{contoh}
\textbf{Contoh 3.} Kriptanalis memegang cipherteks \texttt{HOJKOREGFP} dari Contoh~1.
\begin{itemize}
  \item Dengan kunci \texttt{POYYAEAAZX}, plainteks yang dihasilkan adalah \texttt{SALMONEGGS} (telur ikan salmon).
  \item Dengan kunci \texttt{BXFGBMWCUM}, plainteks yang dihasilkan adalah \texttt{GREENFIELD} (lapangan hijau).
\end{itemize}
Kunci mana yang benar? Kriptanalis tidak punya cara untuk mengetahuinya. Setiap plainteks sepuluh huruf mempunyai tepat satu kunci yang memetakannya ke \texttt{HOJKOREGFP}, yaitu $K = (C - P) \bmod 26$, sehingga cipherteks tidak memberi petunjuk apa pun tentang plainteks yang sebenarnya.
\end{contoh}

\textbf{Pembuktian bahwa cipherteks OTP acak.} Misalkan $C = (P + K) \bmod 26$ dengan $K$ dipilih secara seragam (\textit{uniform}):
\[
\Pr(K = k) = \frac{1}{26}, \qquad k = 0, 1, \ldots, 25.
\]
Untuk setiap nilai plainteks tertentu $P = p$, pemetaan $k \mapsto (p + k) \bmod 26$ adalah pemetaan satu-ke-satu dari semua kemungkinan $K$ ke semua kemungkinan $C$. Akibatnya
\[
\Pr(C = c \mid P = p) = \frac{1}{26} \quad \text{untuk setiap } c,
\]
yaitu $\Pr(C = \texttt{A}) = \Pr(C = \texttt{B}) = \cdots = \Pr(C = \texttt{Z}) = 1/26$, apa pun plainteksnya.

\subsection{Perfect Secrecy}

\begin{figure}[htbp]
  \centering
  \begin{minipage}[c]{0.22\textwidth}
    \centering
    \includegraphics[width=\linewidth]{bab-04/potret-shannon}\\[2pt]
    {\small (a) Claude Shannon}
  \end{minipage}\hfill
  \begin{minipage}[c]{0.72\textwidth}
    \centering
    \includegraphics[width=\linewidth]{bab-04/shannon-1949-makalah}\\[2pt]
    {\small (b) Halaman pertama makalah 1949}
  \end{minipage}
  \caption{Claude Shannon (1916--2001) membuktikan \textit{perfect secrecy} OTP dalam ``Communication Theory of Secrecy Systems'' \cite{shannon1949,munir_slides_itb}.}
  \label{fig:shannon-otp}
\end{figure}

\textit{Perfect secrecy} berarti cipherteks tidak memberikan informasi apa pun tentang plainteks maupun kunci. Pada tahun 1949, Claude Shannon dari Bell Labs membuktikan bahwa OTP memiliki \textit{perfect secrecy} dalam makalah ``Communication Theory of Secrecy Systems'' di \textit{Bell System Technical Journal} 28(4): 656--715 (Gambar~\ref{fig:shannon-otp}) \cite{shannon1949}.

\begin{konsep}
Menurut Shannon, sebuah sistem memiliki \textbf{perfect secrecy} apabila setelah penyerang memperoleh cipherteks, informasinya tentang plainteks tidak bertambah:
\[
\Pr(P = p \mid C = c) = \Pr(P = p) \quad \text{untuk setiap } p \text{ dan } c.
\]
Artinya, peluang suatu plainteks tetap sama sebelum maupun sesudah cipherteks diamati.
\end{konsep}

Pembuktian pada subbagian sebelumnya menunjukkan bahwa $\Pr(C = c \mid P = p)$ tidak bergantung pada $p$; dengan teorema Bayes, hal ini setara dengan definisi di atas. Dari sisi teori informasi, syarat perlunya adalah $H(K) \ge H(P)$: kunci harus mengandung ketidakpastian setidaknya sebesar plainteks (lihat pengayaan \textit{unicity distance} dan \textit{perfect secrecy} di Bab~III).

\subsection{Penyalahgunaan Kunci: Bukan Lagi OTP}

Keamanan OTP runtuh bila salah satu syaratnya dilanggar \cite{munir_slides_itb}:
\begin{itemize}
  \item \textbf{Kunci diambil dari teks panjang} (novel, buku, berita). Ini bukan lagi OTP, melainkan \textit{running-key} Vigen\`{e}re cipher, karena teks tersebut tidak acak. Cipher ini tidak memiliki \textit{perfect secrecy} dan dapat dipecahkan.
  \item \textbf{Kunci dipakai dua kali.} Ini bukan lagi \textit{one-time pad}, melainkan \textit{two-time pad}, dan tidak aman. Jika $C_1 = (P_1 + K) \bmod 26$ dan $C_2 = (P_2 + K) \bmod 26$, maka
  \[
  (C_1 - C_2) \bmod 26 = (P_1 - P_2) \bmod 26.
  \]
  Kunci saling meniadakan, sehingga penyerang memperoleh selisih dua plainteks bermakna yang dapat ditebak dengan statistik bahasa atau kata-kata yang mungkin muncul.
\end{itemize}

\begin{catatan}
Menggunakan kunci berulang, kunci yang dapat diprediksi, atau kunci yang dibangkitkan dengan pembangkit bilangan acak semu menghancurkan keamanan OTP. OTP aman \textbf{hanya} jika kunci acak sejati, sepanjang pesan, dan dipakai sekali.
\end{catatan}

\subsection{Kelemahan OTP dalam Praktik}

Meskipun keamanannya sempurna, OTP jarang dipakai dalam aplikasi praktis, baik komersial maupun lainnya \cite{munir_slides_itb}:
\begin{enumerate}
  \item \textbf{Tidak sangkil.} Panjang kunci sama dengan panjang pesan, sehingga makin panjang pesan makin besar kuncinya. Membangkitkan miliaran karakter yang benar-benar acak membutuhkan komputasi berat, dan kebutuhan kunci menjadi sangat besar untuk komunikasi yang berlangsung terus-menerus.
  \item \textbf{Sulit disinkronkan.} Karena kunci acak sejati, ``tidak mungkin'' pengirim dan penerima membangkitkan kunci yang sama secara serentak. Kunci harus disepakati dan didistribusikan lebih dulu.
  \item \textbf{Butuh saluran kedua yang sangat aman.} Kunci harus dikirim melalui saluran yang berbeda dari saluran pesan, yang biasanya lambat dan mahal (misalnya jalur darat dengan kurir tepercaya yang tidak dikenali), atau melalui saluran yang sama pada waktu yang berbeda.
  \item \textbf{Tidak cocok untuk komunikasi dinamis.} Pada WhatsApp, HTTPS, VPN, penyimpanan awan, surel, dan IoT, data terus berubah, perangkat berganti, pengguna bisa luring, dan ukuran pesan berbeda-beda. Dengan OTP, setiap tambahan data membutuhkan tambahan kunci acak yang belum pernah dipakai, sehingga manajemen kunci menjadi sangat berat.
\end{enumerate}

Karena itu, kriptografi modern memakai pendekatan \textit{computationally secure}: kunci pendek (misalnya 128 bit) yang direntangkan menjadi \textit{keystream} oleh algoritma deterministik (\textit{stream cipher}, Bab~V), dengan keamanan yang bersandar pada sulitnya komputasi, bukan pada teori informasi.

\subsection{OTP dalam Sejarah}

OTP pernah dipakai oleh mata-mata selama Perang Dunia~II. Gambar~\ref{fig:otp-syal} memperlihatkan seorang pejabat \textit{Special Operations Executive} (SOE) Inggris dengan syal sutra bertuliskan kunci OTP yang mudah disembunyikan dan dihancurkan.

\begin{figure}[htbp]
  \centering
  \includegraphics[width=0.45\textwidth]{bab-04/otp-syal-pd2}
  \caption{Pejabat SOE Inggris memperlihatkan syal bertuliskan kunci OTP yang dipakai selama Perang Dunia~II (sumber: John H.\ Reif, \textit{History of Computing}, Duke University) \cite{munir_slides_itb}.}
  \label{fig:otp-syal}
\end{figure}

Pada masa Perang Dingin, OTP dipakai pada \textbf{hotline Moskow--Washington} (1963), saluran komunikasi langsung antara pemimpin Amerika Serikat dan Uni Soviet yang dibangun setelah krisis rudal Kuba untuk menghindari perang nuklir. Hotline ini memakai peralatan OTP \textit{ITT Intelex Teletype} L015 (Gambar~\ref{fig:otp-hotline}). Kunci \textit{one-time pad} disimpan pada pita magnetik yang diterbangkan antara Washington DC dan Moskow \cite{munir_slides_itb}.

\begin{figure}[htbp]
  \centering
  \begin{minipage}[c]{0.32\textwidth}
    \centering
    \includegraphics[width=\linewidth]{bab-04/hotline-teletype}\\[2pt]
    {\small (a) ITT Intelex Teletype L015}
  \end{minipage}\hfill
  \begin{minipage}[c]{0.3\textwidth}
    \centering
    \includegraphics[width=\linewidth]{bab-04/pentagon}\\[2pt]
    {\small (b) Washington (Pentagon)}
  \end{minipage}\hfill
  \begin{minipage}[c]{0.3\textwidth}
    \centering
    \includegraphics[width=\linewidth]{bab-04/kremlin}\\[2pt]
    {\small (c) Moskow (Kremlin)}
  \end{minipage}
  \caption{Hotline Moskow--Washington mengenkripsi pesan dengan OTP memakai ITT Intelex Teletype L015 (Lyndon Baines Johnson Library and Museum) \cite{munir_slides_itb}.}
  \label{fig:otp-hotline}
\end{figure}

\begin{konsep}
``\textit{One-time pads are theoretically unbreakable, but practically very weak. By contrast, conventional ciphers are theoretically breakable, but practically strong.}'' --- Steve Bellovin \cite{munir_slides_itb}

OTP adalah pembanding ideal: aman sempurna secara teori, tetapi lemah dalam praktik karena masalah pembangkitan dan distribusi kunci. Cipher modern menukar keamanan sempurna dengan kepraktisan.
\end{konsep}

